\documentclass[11pt,a4paper]{article}
\usepackage[utf8]{inputenc}
\usepackage[T1]{fontenc}
\usepackage[spanish,es-nodecimaldot]{babel}
\usepackage{amsmath,amssymb,amsthm,bm,geometry}
\usepackage[colorlinks=true,allcolors=blue]{hyperref}
\geometry{margin=25mm}
\newtheorem{proposition}{Proposición}
\newcommand{\PEB}{\operatorname{PEB}}
\title{PEB tridimensional con receptor de respuesta angular generalizada\\\large LED fijo y PD reorientable con $\cos^{m_R}\psi$}
\author{Proyecto Cambridge}
\date{16 de septiembre de 2026}
\begin{document}
\maketitle

\section{Alcance y convenciones}
Esta derivación amplía, sin reemplazar, \path{PEB_derivation.tex}. Se mantiene la arquitectura de un LED fijo y un único PD reorientable. Se conocen la posición y orientación del LED, la potencia, la calibración óptica, el patrón emisor y todas las normales globales del PD. El centro óptico permanece fijo durante el barrido. Se supone propagación LOS y ruido Gaussiano independiente con varianza óptica conocida, constante respecto a posición y señal.

La familia empírica de respuesta del receptor se normaliza respecto a incidencia normal. El nuevo exponente $\nu=m_R>0$ es conocido y calibrado. No es la responsividad espectral en A/W ni un segundo orden de emisión Lambertiana: no se introduce un factor $(m_R+1)$ en la ganancia del detector. La FOV sigue siendo un semiángulo de corte independiente de $m_R$. El valor por defecto es $m_R=1$.

Sean $\mathbf t$ y $\mathbf r$ las posiciones del LED y del PD; el LED tiene normal fija $\mathbf n_t$ y el PD utiliza normales unitarias conocidas $\mathbf n_i$. Definimos:
\begin{equation}
 d=\|\mathbf t-\mathbf r\|,\quad \mathbf u=(\mathbf t-\mathbf r)/d,\quad
 \mathbf q=-\mathbf n_t,\quad c=\mathbf q^{\mathsf T}\mathbf u,\quad s_i=\mathbf n_i^{\mathsf T}\mathbf u.
\end{equation}
La dirección $\mathbf u$ apunta del receptor al transmisor. El orden del emisor y la constante radiométrica son:
\begin{equation}
 m_T=-\frac{\ln2}{\ln(\cos\Phi_{1/2})},\qquad
 C=\frac{P_t(m_T+1)A_{\rm PD}T_s g}{2\pi}.
\end{equation}
En el interior visible $c>0$, $s_i>\cos\Psi$ y $s_i>0$, el modelo y las medias observadas son:
\begin{equation}
 \mu_i=C d^{-2}c^{m_T}s_i^\nu=\beta s_i^\nu,\qquad
 \beta=Cc^{m_T}/d^2,\qquad
 \bar Y_i\sim\mathcal N(\mu_i,\sigma^2/N_i).
\label{eq:model}
\end{equation}
Fuera de aceptación la señal LOS se anula. Todas las muestras siguen consumiendo presupuesto; no se seleccionan observaciones mediante un oráculo de posición.

\section{Gradiente cartesiano}
Las identidades geométricas son:
\begin{equation}
 \nabla d=-\mathbf u,\qquad
 \frac{\partial\mathbf u}{\partial\mathbf r^{\mathsf T}}=-\frac{I-\mathbf u\mathbf u^{\mathsf T}}{d},\qquad
 \nabla c=-\frac{\mathbf q-c\mathbf u}{d},\qquad
 \nabla s_i=-\frac{\mathbf n_i-s_i\mathbf u}{d}.
\end{equation}
La derivada de la amplitud común es:
\begin{equation}
 \nabla\beta=\frac{\beta}{d}\left[(m_T+2)\mathbf u-\frac{m_T}{c}\mathbf q\right].
\end{equation}
La regla del producto aplicada a \eqref{eq:model} proporciona:
\begin{align}
 \nabla\mu_i&=s_i^\nu\nabla\beta+\nu\beta s_i^{\nu-1}\nabla s_i,\\
 &=\frac{Cc^{m_T}}{d^3}\left[-\nu s_i^{\nu-1}\mathbf n_i
 -\frac{m_T}{c}s_i^\nu\mathbf q+(m_T+\nu+2)s_i^\nu\mathbf u\right].
\label{eq:gradient}
\end{align}
Al sustituir $\nu=1$ se recupera exactamente el gradiente anterior con coeficiente radial $(m_T+3)$. Las potencias fraccionarias se evalúan sólo con cosenos positivos en el interior visible, evitando valores complejos fuera de aceptación.

\section{FIM y cota de posición}
Con $J$ formado por las filas $\nabla\mu_i^{\mathsf T}$ y $W=\operatorname{diag}(N_i/\sigma^2)$, la log-verosimilitud y la puntuación son:
\begin{equation}
 \ell=-\frac12\sum_i\left[\ln(2\pi\sigma^2/N_i)+\frac{N_i}{\sigma^2}(\bar Y_i-\mu_i)^2\right],\qquad
 \nabla\ell=J^{\mathsf T}W(\bar{\mathbf Y}-\bm\mu).
\end{equation}
Su segundo momento da la información de Fisher \cite{kay}:
\begin{equation}
 \mathcal I(\mathbf r)=J^{\mathsf T}WJ=\sum_{i\in\mathcal V}\frac{N_i}{\sigma^2}\nabla\mu_i\nabla\mu_i^{\mathsf T},\qquad
 \PEB(\mathbf r)=\sqrt{\operatorname{tr}(\mathcal I^{-1})}.
\label{eq:peb}
\end{equation}
Es una cota local para estimadores regulares localmente insesgados, no un límite universal para estimadores sesgados ni una probabilidad de éxito. La FIM tiene unidades $\mathrm m^{-2}$ y el PEB, metros. En los mapas a una altura dada se conserva la FIM de tres coordenadas; no se supone que la altura es conocida.

Se usa SVD del Jacobiano blanqueado. Si sus valores singulares son $\lambda_j$, se tiene:
\begin{equation}
 \PEB^2=\sum_{j=1}^3\lambda_j^{-2},\qquad
 \mathcal I^{-1}=V\operatorname{diag}(\lambda_j^{-2})V^{\mathsf T}.
\end{equation}
La implementación devuelve infinito si el rango numérico es menor que tres, con tolerancia relativa $10^{-10}$. Conserva la exclusión de fronteras mediante \texttt{NaN} y tolerancia $10^{-12}$ en coseno. Con FOV menor que $90^\circ$ hay un salto de potencia en el corte. En incidencia rasante a $90^\circ$, la suavidad depende de $m_R$; para órdenes suficientemente altos puede existir una extensión diferenciable. Excluir también esos casos límite es una convención conservadora del código, no afirmar que toda frontera sea siempre matemáticamente no regular.

\section{Identificabilidad y transformación de una sola etapa}
En una región visible fija definimos el vector óptico:
\begin{equation}
 \mathbf v=\beta^{1/\nu}\mathbf u,\quad a=\|\mathbf v\|=\beta^{1/\nu},\qquad
 \mu_i^{1/\nu}=\mathbf n_i^{\mathsf T}\mathbf v.
\label{eq:rootlinear}
\end{equation}
Si las normales visibles tienen rango tres, las potencias sin ruido determinan $\mathbf v$. La conversión a posición es algebraica:
\begin{equation}
 \mathbf u=\mathbf v/a,\qquad \beta=a^\nu,\qquad
 d=\sqrt{Cc^{m_T}/a^\nu},\qquad \mathbf r=\mathbf t-d\mathbf u.
\label{eq:inverse}
\end{equation}
No hace falta una adquisición de distancia adicional ni dos optimizaciones independientes.

El Jacobiano de esa transformación es:
\begin{equation}
 B_\nu=\frac{\partial\mathbf v}{\partial\mathbf r^{\mathsf T}}
 =-\frac{a}{d}\left[I+\mathbf u\left(\frac{m_T}{\nu c}\mathbf q^{\mathsf T}
 -\frac{m_T+\nu+2}{\nu}\mathbf u^{\mathsf T}\right)\right].
\end{equation}
Por el lema del determinante:
\begin{equation}
 \det B_\nu=\left(-\frac ad\right)^3\left[1+\frac{m_T}{\nu}-\frac{m_T+\nu+2}{\nu}\right]
 =\frac{2a^3}{\nu d^3}>0.
\end{equation}
Como:
\begin{equation}
 J_{\mathcal V}=\nu\operatorname{diag}\left[(H_{\mathcal V}\mathbf v)^{\nu-1}\right]H_{\mathcal V}B_\nu,
\end{equation}
los factores laterales son invertibles en el interior visible y el rango es exactamente el de $H_{\mathcal V}$. Para cualquier $m_R>0$ conocido, siguen siendo necesarias y suficientes tres normales visibles linealmente independientes para identificación local 3D. El orden cambia las magnitudes, la ponderación de filas y el condicionamiento, aunque no el rango matemático. Los órdenes extremos pueden afectar al rango numérico efectivo. El caso $m_R=0$ es diferente: la respuesta es constante dentro de FOV y la diversidad angular regular desaparece; por ello no está incluido en la familia implementada.

\section{Caso axial y leyes de escala}
En el eje de un LED vertical, con cono uniforme, $K\geq3$, $N_i=N$ y todas las orientaciones visibles:
\begin{equation}
 \mathcal I=\frac{NKC^2}{\sigma^2d^6}\operatorname{diag}\left(
 \frac{\nu^2\sin^2\theta\cos^{2\nu-2}\theta}{2},
 \frac{\nu^2\sin^2\theta\cos^{2\nu-2}\theta}{2},
 4\cos^{2\nu}\theta\right).
\end{equation}
De ahí:
\begin{equation}
 \PEB^2=\frac{\sigma^2d^6}{NKC^2}\left[
 \frac4{\nu^2\sin^2\theta\cos^{2\nu-2}\theta}+\frac1{4\cos^{2\nu}\theta}\right].
\label{eq:axial}
\end{equation}
Para $\nu=1$ se recupera el caso anterior y su óptimo axial de $\arctan2$. Ese ángulo no es universal para otros órdenes. En el interior visible se mantienen las leyes $\PEB\propto1/A_{\rm PD}$, $1/P_t$ y $1/\sqrt N$ bajo ruido óptico constante. No debe suponerse que aumentar $m_R$ siempre mejora: aumenta la pendiente angular pero también atenúa la señal a incidencias oblicuas.

\section{Patrón del LED conocido, pero no exactamente Lambertiano}
Sea $f_T(\mathbf u)>0$ una respuesta angular calibrada. Basta sustituir $c^{m_T}$ por $f_T$ y definir $\beta=Cf_T/d^2$. La derivada general, usando el gradiente tangente $\nabla_S f_T=(I-\mathbf u\mathbf u^{\mathsf T})\nabla_{\mathbf u}f_T$, es:
\begin{equation}
 \nabla\mu_i=\frac{\beta}{d}\left[-\nu s_i^{\nu-1}\mathbf n_i
 +(\nu+2)s_i^\nu\mathbf u-s_i^\nu\nabla_S\ln f_T\right].
\end{equation}
La transformación \eqref{eq:rootlinear} y su determinante $2a^3/(\nu d^3)$ siguen siendo válidos. El patrón emisor entra en la amplitud y en la conversión a distancia; los cocientes entre orientaciones cancelan ese factor común. Esto no hace innecesaria su calibración para posición 3D.

Para ensayos controlados se implementa una perturbación azimutal explícita, no un ajuste universal de LEDs reales:
\begin{equation}
 f_T(\mathbf u)=c^{m_T}\left[1+\epsilon\left((\mathbf e_1^{\mathsf T}\mathbf u)^2-(\mathbf e_2^{\mathsf T}\mathbf u)^2\right)\right],\qquad |\epsilon|<1.
\end{equation}
$\mathbf e_1,\mathbf e_2$ son ejes transversales conocidos del LED. La corrección es positiva y su promedio azimutal es cero, por lo que mantiene la normalización de potencia total del patrón Lambertiano. Por defecto $\epsilon=0$. El código utiliza el gradiente completo de este patrón cuando se activa, no el gradiente Lambertiano aplicado a una media distinta.

\section{Supuestos de calibración y verificación}
Si $C$ es desconocido, la transformación $d\mapsto kd$, $C\mapsto k^2C$ sigue conservando todas las medias. Si $m_R$ es incierto, debe tratarse como parámetro auxiliar; su derivada es $\partial\mu_i/\partial m_R=\mu_i\ln s_i$. La cota anterior presupone que ya se ha calibrado. A incidencia normal en un cono uniforme, ese parámetro puede confundirse con amplitud, por lo que no debe inferirse de un único barrido simétrico sin análisis de rango.

Se comprueban: coincidencia de medias, gradientes y PEB con la fórmula anterior cuando $m_R=1$; diferencias centrales para órdenes fraccionarios y enteros; expresión axial \eqref{eq:axial}; rango y máscara FOV; ausencia de valores complejos fuera de aceptación; y derivadas del patrón emisor no Lambertiano. Los scripts de diseño conservan sus gráficas y añaden únicamente el parámetro físico y su registro. Los MAT anteriores sin ese campo se interpretan como $m_R=1$.

\bibliographystyle{IEEEtran}
\bibliography{PEB_references}
\end{document}
